\section{Dataset sử dụng trong demo}
\label{sec:dataset}

\subsection{Hai benchmark công khai}

Demo dùng hai benchmark thật, đại diện cho hai \emph{construct} khác nhau:

\begin{itemize}
  \item \textbf{MMLU}~\cite{hendrycks2021mmlu} (Massive Multitask Language
    Understanding) --- đo \emph{kiến thức/hiểu biết} qua trắc nghiệm 4 đáp án trải
    trên 57 môn (từ toán, luật, y đến đạo đức). Mỗi câu có đúng một đáp án
    \texttt{A/B/C/D}.
  \item \textbf{GSM8K}~\cite{cobbe2021gsm8k} (Grade-School Math 8K) --- đo
    \emph{suy luận toán} qua bài toán đố tiểu học; đáp án là một \emph{số} cuối
    cùng, lời giải nhiều bước.
\end{itemize}

Việc chọn \emph{hai} construct là có chủ đích: nó cho phép kiểm tra luận điểm
trung tâm của tutorial --- \emph{một benchmark không xếp hạng được model};
liệu thứ hạng trên MMLU và GSM8K có khớp nhau không (Phần~\ref{sec:results}).

\subsection{Lấy mẫu cố định theo seed}

Để demo tái lập được trên máy free (T4), nhóm lấy \emph{tập con} cố định
thay vì toàn bộ benchmark:

\begin{table}[t]
\centering
\caption{Tham số lấy tập con benchmark cho demo (cố định, ghi trong
\texttt{config.py}).}
\label{tab:subset}
\small
\begin{tabular}{lcc}
\hline
\textbf{Tham số} & \textbf{Ký hiệu} & \textbf{Giá trị} \\
\hline
Số câu MMLU & \texttt{N\_MMLU} & 150 \\
Số câu GSM8K & \texttt{N\_GSM8K} & 80 \\
Số câu MMLU được sinh nhiễu & \texttt{N\_PERTURB} & 50 \\
Seed cố định & \texttt{SEED} & 0 \\
\hline
\end{tabular}
\end{table}

MMLU được lấy \emph{phân tầng} (stratified) trên cả 57 môn rồi bù thêm cho đủ 150
câu, tránh việc một vài môn lấn át. Mọi item và thứ tự được ghi ra
\texttt{data/*.jsonl}, nên \emph{bất kỳ ai chạy lại đều nhận đúng cùng tập câu, cùng
thứ tự}. Đây là yêu cầu \emph{reproducibility} của BetterBench được thực thi ở mức
dữ liệu.

\subsection{Bản sinh đôi nhiễu (perturbation)}

Lấy cảm hứng từ tinh thần SuperGLUE (ra đề khó hơn khi đề cũ bão hoà), mỗi câu MMLU
trong tập 50 câu có một \emph{bản sinh đôi giữ nhãn}: nội dung và đáp án đúng không
đổi, chỉ thay đổi bề mặt theo một trong hai phép:

\begin{itemize}
  \item \texttt{option\_shuffle}: đảo vị trí các lựa chọn (đáp án đúng đổi chữ cái
    nhưng vẫn là cùng nội dung) --- bắt mô hình ``đoán theo vị trí'' sẽ lộ.
  \item \texttt{distractor}: chèn thêm lựa chọn ``None of the above'' --- tăng
    nhiễu mà không đổi đáp án đúng.
\end{itemize}

Nếu accuracy tụt mạnh giữa bản gốc và bản nhiễu, điểm gốc một phần dựa vào
\emph{format/vị trí} chứ không thuần kiến thức (xem Phần~\ref{sec:results}).

\subsection{Phạm vi và giới hạn của dataset}

Tập con nhỏ (150 + 80) khiến khoảng tin cậy rộng; mục tiêu là \emph{minh hoạ vấn
đề đo lường}, không phải xếp hạng tuyệt đối. Quan trọng hơn, cả MMLU (2021) và
GSM8K (2021) gần như chắc đã nằm trong dữ liệu huấn luyện của các model hiện đại
--- một nguồn \emph{nhiễm dữ liệu} mà demo \emph{cố ý} giữ lại để phơi bày, chứ
không che giấu (Phần~\ref{sec:results}).
